\section{Related literature}
\label{sec:literature}
Our study builds on three strands of work: incentives for exercise and the habits they may create, commitment devices, and reminders and plans.

\subsection{Incentives and habits}
A first strand pays people to exercise and asks whether a habit survives once the payments stop. \citet{charness2009} find lasting increases in attendance after a month of paid gym visits, driven by people who had not been regular gym users before. \citet{acland2015} replicate this design with students and staff of a large public university and find that the habit formed during the incentive month largely decayed over the winter break. \citet{royer2015} pay employees of a Fortune 500 company for visits to the onsite gym and find strong responses while the payments last but little persistence, unless the incentive is followed by the offer of a commitment contract. Among new members of a private gym, \citet{carrera2018incentives} find that modest threshold incentives had small effects on attendance, and that the effects on average visits were largest among members who already exercised regularly.

\subsection{Commitment devices}
A second strand studies commitment devices directly. Theory predicts a demand for commitment from people who are aware of their present bias \citep{odonoghue1999}, and \citet{bryan2010} review commitment devices in saving, health and other domains. \citet{augenblick2015} find considerable present bias in the allocation of effort but very little in choices over money, and 59 percent of their subjects chose a commitment device when it was offered at a price of zero. In a data-entry firm in India, workers chose contracts that penalized them for low output \citep{kaur2015}. Deposit contracts, in which participants put their own money at risk, helped obese adults lose weight in a randomized trial \citep{volpp2008}.

Commitment contracts for exercise have been studied in workplaces and online. In \citet{royer2015}, 22 percent of the employees offered a self-funded commitment contract signed one. Those who committed placed 58 dollars at stake on average, and 63 percent of them met their target. \citet{bhattacharya2015} randomize the default length of exercise contracts on the stickK.com platform. Longer defaults led users to choose longer contracts but also made them less likely to sign a contract at all. Only 22.7 percent of the signed contracts put money at stake.

Commitment does not have to involve money. \citet{milkman2014} let students, faculty and staff of a university listen to tempting audiobooks only at the gym. Attendance rose until the Thanksgiving break, and the program helped the busiest participants most. When offered the chance to pay for a month of gym-only access to a tempting audiobook, participants in that study were willing to pay \$9.61 on average.

Recent work questions whether the take-up of contracts reveals awareness of self-control problems. In an experiment with 1,248 members of a fitness center, \citet{carrera2022} find high take-up of contracts that pay only if attendance reaches a threshold (64, 49 and 32 percent for thresholds of 8, 12 and 16 visits in four weeks), but an information treatment that made members more aware of their overoptimism reduced demand for contracts. They estimate an average present-bias parameter of 0.84 and a perceived parameter of 0.55, so that members are more pessimistic about their future self-control than their behavior warrants.

\subsection{Reminders and plans}
A third strand studies reminders and plans. Reminders help clients of commitment savings accounts reach their savings goals \citep{karlan2016}, and prompts to write down the date and time of a planned flu shot raise vaccination rates \citep{milkman2011}. For repeated behavior the evidence is weaker. \citet{carrera2018plans} asked members of a private gym to plan their visits for the next two weeks and find a precisely estimated null effect on attendance; total planned visits were three times actual visits (7.04 against 2.33). In the megastudy of \citet{milkman2021}, which tested 54 four-week programs with 61,293 members of a gym chain, a basic condition combining planning prompts, reminders and small rewards raised weekly visits by 0.14, or 9 percent. The authors report unadjusted p-values for their 53 comparisons together with an estimate of the false discovery rate.

We add a direct test of whether reminders make commitment contracts more effective, with contracts that participants fund themselves and with randomization across departments rather than individuals.
